% =====================================================================
%  CHƯƠNG 2. TỔNG QUAN CÔNG TRÌNH LIÊN QUAN       | Phụ trách: Cả 2
%  Ngân sách: mục tiêu 8 trang, trần 9 trang
%  Nguồn: docs/research/2026-07-18-yolo-architecture.md            (2.1),
%         docs/research/khao-sat-he-thong-alpr.md                  (2.2, 2.3),
%         docs/research/2026-07-19-similar-parking-systems.md      (2.4, 2.5)
% =====================================================================
\chapter{Tổng Quan Công Trình Liên Quan}
\label{ch:lienquan}

\section{Phát Hiện Đối Tượng và Họ Kiến Trúc YOLO}

\subsection{Đối Lập Một Giai Đoạn và Hai Giai Đoạn}

Các phương pháp phát hiện đối tượng (object detection) hai giai đoạn (two-stage) sinh
tập vùng đề xuất (region proposal) rồi phân loại và tinh chỉnh tọa độ từng vùng. Faster
R-CNN vẫn xử lý tuần tự theo vùng đề xuất nên chỉ đạt khoảng năm FPS trên GPU cùng
thời~\autocite{ren2015fasterrcnn}. Ngược lại, họ YOLO (You Only Look
Once) đưa phát hiện về một bài toán regression duy nhất: ảnh chia lưới, mạng dự đoán
trực tiếp bounding box và lớp trong một lần forward pass. YOLOv1 đạt 45 FPS với mAP gần
gấp đôi các detector real-time cùng thời, đánh đổi một phần độ chính xác định vị để lấy
tốc độ~\autocite{redmon2016yolo}. Đánh đổi này phù hợp ràng buộc thời gian thực trên CPU
không GPU: một detector hai giai đoạn chạy hai lần liên tiếp cho phát hiện xe rồi biển số
sẽ cộng dồn độ trễ tới mức khó đạt mục tiêu dưới hai giây.

\subsection{Kiến Trúc YOLOv8}
\label{sec:yolov8}

YOLOv8 đánh dấu bước chuyển kiến trúc quan trọng: bỏ hoàn toàn anchor box (anchor-free),
tách head detection thành hai nhánh song song (decoupled head) cho box regression và
class score, và dùng Task-Aligned Assigner khi gán nhãn huấn luyện. Backbone dùng block
C2f thay cho C3 của YOLOv5 để tái sử dụng đặc trưng phong phú hơn; neck PAN-FPN fuse
feature ba tỷ lệ. Mỗi anchor point regress bốn khoảng cách cạnh qua Distribution Focal
Loss (DFL) với 16 bin xác suất rời rạc thay vì một scalar, cải thiện định vị vật thể nhỏ
và mờ; điểm class sau sigmoid chính là confidence, không có nhánh objectness riêng. Sau
giải mã DFL, Non-Maximum Suppression (NMS) lọc các detection trùng lặp, là bước hậu xử lý
tuần tự nằm ngoài graph tensor, thêm độ trễ và phức tạp hóa xuất
ONNX~\autocite{jocher2023yolov8,terven2023comprehensive}.

\subsection{YOLO26: Thiết Kế Hướng Đến Edge và Xuất Mô Hình}
\label{sec:yolo26}

YOLO26, ra mắt tháng 1/2026 với paper arXiv công bố tháng 6/2026, được Ultralytics định
vị là thế hệ hiện tại cho deploy CPU và edge. Bốn thay đổi kiến trúc chính so với YOLOv8
được xác nhận từ abstract arXiv và docs chính
thức~\autocite{jocher2026yolo26,chakrabarty2026yolo26nmsfree}:

\begin{itemize}
  \item \textbf{NMS-free end-to-end:} dual-head cho phép train theo cách gán one-to-many thông thường nhưng export và chạy bằng head one-to-one xuất detection cuối trực tiếp, không cần bước NMS, đơn giản hóa graph ONNX và giảm latency deploy.
  \item \textbf{Bỏ DFL:} head không còn softmax 16-bin cho mỗi cạnh box, nên head nhẹ hơn và graph export INT8 đơn giản hơn vì loại op Softmax vốn khó lượng tử hóa trên NPU và DSP~\autocite{chakrabarty2026yolo26nmsfree}.
  \item \textbf{MuSGD:} optimizer lai SGD và Muon lấy cảm hứng từ huấn luyện mô hình ngôn ngữ lớn, cải thiện hội tụ.
  \item \textbf{Progressive Loss và STAL:} Progressive Loss dịch trọng tâm supervision về head one-to-one; STAL (Small-Target-Aware Label Assignment) đảm bảo độ phủ positive cho vật thể nhỏ, nhắm trực tiếp vào bài toán biển số nhỏ trong khung ảnh bãi đỗ.
\end{itemize}

Trên COCO val với input 640, YOLO26n đạt 40,9 mAP ở 38,9 ms CPU ONNX (Intel Xeon), còn
YOLOv8n đạt 37,3 mAP ở 80,4 ms cùng điều kiện.

Đề tài so sánh trực tiếp YOLOv8n và YOLO26n trên dữ liệu biển số Việt Nam ở
Chương~\ref{ch:phathien}.

\section{Nhận Dạng Biển Số Xe Tự Động}

\subsection{Kiến Trúc Ba Tầng trong ALPR}

Nhận dạng biển số xe tự động (Automatic License Plate Recognition, ALPR) thường gồm ba
tầng: phát hiện vùng biển, chuẩn hóa và tách dòng, rồi đọc ký tự bằng OCR hoặc theo hướng
OCR-as-detection (dùng detector YOLO phát hiện từng ký tự). Al-batat và cộng sự công bố
một pipeline ALPR mã nguồn mở đánh giá trên năm bộ dữ liệu nhiều khu vực (Mỹ, EU, Đài
Loan, Brazil), đạt 90,3\% ở cấp ký tự trung bình; khi đo từng tầng riêng, phát hiện xe
đạt 99,71\% precision và nhận dạng ký tự tách riêng đạt 99,68\%, cho thấy khoảng cách lớn
giữa kết quả từng tầng và kết quả toàn
pipeline~\autocite{albatat2022endtoend,redalb2022alprcode}.

Bằng chứng định lượng rõ nhất cho lợi ích kiến trúc đa tầng là so sánh trực tiếp của
Safran và cộng sự: đa tầng đạt 96,1\% trong khi single-stage chỉ 83,9\% trên cùng dữ liệu
và điều kiện, chênh hơn 12 điểm phần trăm~\autocite{safran2024multistage}. Hệ thống này
còn tích hợp web dashboard trực quan hóa trạng thái chiếm chỗ và luồng xe thời gian thực,
một trong số ít công trình kết hợp cả ALPR và dashboard.

Moussaoui và cộng sự dùng YOLOv8 kết hợp EasyOCR với tiền xử lý k-means, thresholding và
morphology, đạt 99\% phát hiện và 98\% ký tự trên tập nhỏ 270
ảnh~\autocite{moussaoui2024yolov8ocr}.

\subsection{Bộ Dữ Liệu Quốc Tế và Giới Hạn với Biển Số Việt Nam}

Các benchmark quốc tế cung cấp nền tảng so sánh nhưng không phản ánh đặc thù biển số Việt
Nam. CCPD (Chinese City Parking Dataset) gồm khoảng 290 nghìn ảnh chụp tại bãi xe Bắc
Kinh với đa dạng điều kiện ánh sáng, góc nghiêng và thời tiết, là benchmark bãi xe lớn
nhất tìm được~\autocite{xu2018ccpd}. UFPR-ALPR gồm 4500 ảnh từ 150 phương tiện tại Brazil
ở độ phân giải 1920×1080, xe và camera đều chuyển động, có gán nhãn vị trí từng ký
tự~\autocite{laroca2018ufpr}. Cả hai chỉ có biển một hàng của nước ngoài, không
có biển hai hàng và ký tự ``Đ'' của Việt Nam, nên không thể so sánh số tuyệt đối với đề
tài.

\section{Nghiên Cứu trên Biển Số Việt Nam}

\subsection{Các Công Trình Nhận Dạng Biển Số Việt Nam}

Trong nước, một số công trình xử lý trực tiếp bối cảnh gần với đề tài. Tran và Bui triển
khai SSD-MobileNetV2 phát hiện vùng biển kết hợp YOLOv8-nano nhận dạng ký tự trên
Raspberry Pi 4 Model B với camera 8 MP, đạt độ chính xác trung bình 95,68\% với thời gian
0,478 giây mỗi ảnh~\autocite{tran2025vietnamlpr}. Đây là một trong số ít công trình trong
nước đo cả độ chính xác lẫn thời gian trên phần cứng biên (số liệu thứ cấp, chưa tiếp cận
được bản đầy đủ), nhưng dừng ở khâu nhận dạng, không có quản lý phiên hay tính phí.

Dang và cộng sự đề xuất YOLO phát hiện xe, WPOD-NET trích vùng biển (xử lý tốt biển xiên
góc), rồi CRNN huấn luyện đồng thời với CTC loss và attention để đọc ký tự, đạt Word Error
Rate 0,014 trên dữ liệu tự thu tại bãi đỗ trong nhà ở Việt
Nam~\autocite{dang2024vietnamcrnn}. Ngữ cảnh sát bài toán của đề tài, nhưng cỡ tập test và
protocol đo không được công bố nên khó đánh giá độ tin cậy.

Với biển xe máy, Le và cộng sự đề xuất YOLOv8 ba tầng (phát hiện xe máy, phát hiện biển
trong bbox xe máy, nhận dạng biển), đạt mAP 93\% sau 300 epoch, cung cấp baseline trực
tiếp cho biển hai hàng, loại khó nhất trong bối cảnh đề tài~\autocite{le2023motorcycleplate}.
Các hướng khác gồm nhận dạng đa góc nhìn trên bộ PTITPlates~\autocite{trananh2023multiangle},
cuộc thi nhận dạng biển xe máy UIT MAPR 2018 với 3000 ảnh chụp tại bãi giữ
xe~\autocite{mapr2018challenge}, và dự án mã nguồn mở mrzaizai2k dùng KNN cùng OpenCV cho
cả biển một và hai hàng~\autocite{mrzaizai2k2025vnplate}.

\subsection{Đặc Điểm Chung và Khoảng Hở của Các Công Trình Việt Nam}

Tổng hợp cho thấy hai điểm chung: phần lớn dừng ở khâu nhận dạng biển số, không mở rộng
sang quản lý phiên hay tính phí; và không công trình nào đo hiệu năng trên Raspberry Pi 5
hoặc phần cứng tương đương trong điều kiện thực tế cổng bãi xe.

\section{Hệ Thống Quản Lý Bãi Giữ Xe}

\subsection{Các Hệ Thống Đại Diện}

Pradhan và cộng sự xây dựng bãi xe IoT trên Raspberry Pi 4 quản lý bốn camera, chạy YOLO
cùng Tesseract OCR với tiền xử lý OpenCV, tích hợp cảm biến IR và ESP32 cho mỗi ô đỗ, ghi
thời gian vào ra, tính phí động theo giờ cao và thấp điểm, theo dõi trạng thái chiếm chỗ.
Trên 100 test case: khớp biển tổng thể 88,9\%; theo điều kiện 95\% ban ngày, 90\% ánh
sáng yếu, 93\% góc 45°; sai số theo dõi chiếm chỗ dưới 5\% so với đếm thủ
công~\autocite{pradhan2025iotparking}. Đây là công trình tích hợp trọn vẹn nhất về thiết bị
biên, ALPR, phiên và tính phí, nhưng mỗi ô đỗ cần một ESP32 và cảm biến IR nên chi phí
phần cứng cao.

Ammar và cộng sự triển khai pipeline năm tầng tại cổng campus trên Jetson Xavier AGX:
YOLOv4 phát hiện xe và biển từ RTSP stream, Xception phân loại 196 đời xe, YOLOv4 nhận
dạng ký tự, DeepSORT theo dõi đa khung hình kèm voting, tối ưu bằng TensorRT, đạt 17,1 FPS
trung bình. Voting qua nhiều frame tăng độ chính xác biển đầy đủ từ 29\% (một frame) lên
69\% (tối đa 35 frame); độ phân giải camera quyết định kết quả: 80\% ở 1080p, 72\% ở
720p, chỉ 23\% ở 480p. Khoảng cách giữa validation cấp ký tự (99,8\%) và thực tế toàn
pipeline (74\%) cho thấy cần báo cáo cả hai cấp độ đo~\autocite{ammar2023multistage}.

Rani và cộng sự xây dựng hệ thống có quản lý vào ra và tính phí: module vào ghi biển và
timestamp, module ra tính phí theo thời lượng và sinh QR code động cho thanh toán không
tiếp xúc, đạt 97,11\% LPR khi không tính khoảng trắng và 91,91\% khi tính cả khoảng
trắng, minh họa mức phụ thuộc của kết quả vào quy tắc chuẩn hóa
chuỗi~\autocite{rani2024ips}. Arukonda và cộng sự dùng YOLOv8 kết hợp EasyOCR và Tesseract
với bước validate định dạng biển theo quy tắc quốc gia, đạt 98,5\% detect, latency 50 ms
mỗi frame (phần cứng không nêu), nhưng backend chỉ ghi log Excel chứ không có cơ sở dữ
liệu thực~\autocite{arukonda2026smartparking}.

\subsection{So Sánh Các Hệ Thống Bãi Giữ Xe}

Bảng~\ref{tab:parking-compare} so sánh các hệ thống trên theo chức năng và nền tảng triển
khai.

\begin{table}[htbp]
\centering
\caption{So sánh các hệ thống quản lý bãi giữ xe tiêu biểu.}

\footnotesize Nguồn: tổng hợp từ \autocite{pradhan2025iotparking,ammar2023multistage,safran2024multistage,rani2024ips,arukonda2026smartparking,platerecognizer2026alpr,openalpr2016benchmarks}.
\label{tab:parking-compare}
\begin{tabular}{p{3.0cm} p{2.4cm} p{1.9cm} p{1.7cm} p{1.4cm} p{1.5cm}}
\hline
\textbf{Hệ thống} & \textbf{Nền tảng} & \textbf{Phát hiện biển} & \textbf{Quản lý session} & \textbf{Dashboard} & \textbf{Màu biển} \\
\hline
Pradhan 2025~\autocite{pradhan2025iotparking}
  & Raspberry Pi 4
  & 88,9\% (n=100)
  & Có (tính phí động)
  & Không
  & Không \\
Ammar 2023~\autocite{ammar2023multistage}
  & Jetson Xavier AGX
  & 74\% (video)
  & Không
  & Không
  & Không \\
Safran 2024~\autocite{safran2024multistage}
  & Server
  & 96,1\%
  & Không
  & Có (web)
  & Không \\
Rani 2024~\autocite{rani2024ips}
  & PC
  & 97,1\%
  & Có (QR thanh toán)
  & Không
  & Không \\
Arukonda 2026~\autocite{arukonda2026smartparking}
  & Không rõ
  & 98,5\%
  & Không
  & Không
  & Không \\
Plate Recognizer~\autocite{platerecognizer2026alpr}
  & SDK on-prem (Pi, Jetson)
  & Không công bố
  & Có (ParkPow)
  & Có (ParkPow)
  & Màu xe \\
Rekor Scout~\autocite{openalpr2016benchmarks}
  & Server và edge agent
  & Không công bố
  & Không
  & Có (cloud)
  & Màu xe \\
\textbf{Đề tài này}
  & \textbf{Raspberry Pi 5}
  & \textbf{Mục tiêu >90\%}
  & \textbf{Có (CSDL, tính phí)}
  & \textbf{Có (React)}
  & \textbf{Màu nền biển} \\
\hline
\end{tabular}
\end{table}

Ngoài đề tài này, không hệ thống nào nhận màu nền biển; Plate Recognizer và Rekor Scout chỉ
nhận màu xe. Về deploy, chỉ Pradhan 2025 và đề tài này chạy trên phần cứng edge
Raspberry Pi; các hệ thống còn lại dùng máy chủ riêng hoặc cloud. Về nghiệp vụ, không
công trình học thuật nào khảo sát tích hợp đầy đủ ALPR, session với đối chiếu chống gian
lận, tính phí và dashboard trên một thiết bị biên duy nhất.

\section{Trí Tuệ Nhân Tạo trên Thiết Bị Biên}

\subsection{Nén Mô Hình và Lượng Tử Hóa}

Triển khai suy luận trên thiết bị biên đòi hỏi nén mô hình (model compression) cho phù
hợp bộ nhớ và năng lực tính toán hạn chế. Lượng tử hóa INT8 (quantization) là kỹ thuật
phổ biến nhất: biểu diễn trọng số và kích hoạt từ dấu phẩy động 32-bit (FP32) về số
nguyên 8-bit, giảm bộ nhớ và tăng tốc tính toán nhờ phép nhân nguyên hiệu quả, đổi lại
một phần nhỏ độ chính xác.

Sonnara và cộng sự với kiến trúc Light-Edge (backbone ResNet-18 cùng FPN, head
anchor-free, head nhận dạng CTC) cho thấy lượng tử hóa INT8 qua TensorRT tăng
throughput 4,6 lần từ 3,1 FPS lên 14,2 FPS trên Jetson Nano với input 1280×720, chỉ mất
0,4 điểm phần trăm mAP (từ 90,6\% xuống 90,2\%) nhờ giữ lớp đầu và lớp cuối ở FP16 (lượng
tử hóa hỗn hợp)~\autocite{sonnara2025lightedge}. Jetson Nano và Pi 5 khác lớp phần cứng nên con số
này không chuyển trực tiếp sang Pi 5 được.

Chakrabarty chỉ ra lớp Softmax của DFL khó lượng tử hóa trên NPU và DSP, nên việc YOLO26
bỏ DFL (Mục~\ref{sec:yolo26}) thuận lợi hơn cho INT8~\autocite{chakrabarty2026yolo26nmsfree};
lợi ích này trên Pi 5 với ONNX Runtime vẫn cần đo thực tế.

\subsection{Raspberry Pi 5 và Benchmark Thiết Bị Biên}
\label{sec:pi5bench}

Raspberry Pi 5 với CPU ARM Cortex-A76, ra mắt cuối 2023, là thế hệ tiến bộ đáng kể so với
Pi 4. Hướng dẫn chính thức của Ultralytics báo cáo YOLO26n trên Pi 5 ở 130,33 ms qua ONNX
Runtime và 67,69 ms qua NCNN FP32 với input 640×640~\autocite{ultralytics2026rpiguide}.
Với cascade hai lần detect (xe rồi biển số trên crop), tổng thời gian detection ước tính
khoảng 260 ms qua ONNX, để lại hơn 1,5 giây trong ngân sách 2 giây cho OCR, phân loại màu
và vào ra, nên detection không phải điểm nghẽn chính trên Pi 5 ở FP32. Benchmark đa SBC
xác nhận biến thể nano và small cho hiệu năng tốt nhất trên SBC và các mô hình YOLO
thường chiếm 4 đến 12\% RAM khả dụng trên Pi 4/5~\autocite{mdpi2026sbcbenchmark}, là cơ sở
để đề tài chọn Pi 5 làm nền tảng biên minh họa với biến thể nano cho cả hai detector.

\section{Khoảng Trống Nghiên Cứu và Định Vị Đề Tài}

\subsection{Hai Khoảng Trống Chính}

\textbf{Khoảng trống 1: Chưa hệ thống nào tích hợp trọn vẹn trên một thiết bị biên.}
Bảng~\ref{tab:parking-compare} cho thấy không công trình học thuật hay sản phẩm thương mại
nào kết hợp đồng thời ALPR, backend nghiệp vụ (cơ sở dữ liệu, session vào ra, đối chiếu
chống gian lận, tính phí) và dashboard quản lý trên một thiết bị biên duy nhất. Pradhan
2025 gần nhất về tích hợp nhưng vẫn dùng máy chủ ngoài cho SQL Server và không có dashboard
web; Safran 2024 có dashboard nhưng chạy trên server. Vận hành mọi thành phần trên một Pi
5 buộc suy luận, backend và cơ sở dữ liệu chia sẻ cùng tài
nguyên~\autocite{pradhan2025iotparking,ammar2023multistage}.

\textbf{Khoảng trống 2: Màu nền biển số chưa được khai thác như tín hiệu phân loại.} Tại
Việt Nam, màu nền biển mang giá trị pháp lý, phân biệt xe cá nhân, xe kinh doanh, xe cơ
quan nhà nước và xe quân sự (Mục~\ref{sec:maunenbien}). Thông tin này cho phép phân nhóm phương
tiện ngay từ tầng thị giác mà không cần tra cứu cơ sở dữ liệu ngoài, hỗ trợ áp biểu phí
theo nhóm và phát hiện bất thường. Tuy nhiên không hệ thống nào trong khảo sát dùng tín
hiệu này.

\subsection{Định Vị Đề Tài}

Đề tài lấp hai khoảng trống này bằng pipeline nhận diện năm thành phần
(Mục~\ref{sec:donggop}) gắn với backend nghiệp vụ và dashboard, toàn bộ vận hành được trên
một Raspberry Pi 5.
